\documentclass[11pt,a4paper]{scrartcl}

% ==================================================
% Pakete
% ==================================================
\usepackage[T1]{fontenc}
\usepackage[ngerman]{babel}
\usepackage{lmodern}
\usepackage{mathtools}
\usepackage{booktabs} %professionelle Formatierung für Tabellen
\usepackage{siunitx}

% Deutsche Zahlenformatierung
\sisetup{
	locale = DE,
	per-mode = symbol
}

\begin{document}
	
	\section{Physikalische Größen und Einheiten
		mit \texttt{siunitx}}
	
	% ==================================================
	\subsection{Zahlen formatieren mit \texttt{\textbackslash num}}
	% ==================================================
	
	\begin{itemize}
		\item Große Zahl: \num{1234567.89}
		\item Kleine Zahl: \num{0.000012345}
		\item Negative Zahl: \num{-1234.567}
		\item Wissenschaftliche Notation: \num{1.23e6}
		\item Gerundete Zahl:
		\num[
		round-mode=places,
		round-precision=2
		]{3.14159265}
	\end{itemize}
	
	% ==================================================
	\subsection{Physikalische Größen mit \texttt{\textbackslash qty}}
	% ==================================================
	
	\begin{itemize}
		\item Länge: \qty{12.5}{\metre}
		\item Masse: \qty{75.3}{\kilogram}
		\item Zeit: \qty{120}{\second}
		\item Temperatur: \qty{22.5}{\celsius}
		\item Geschwindigkeit:
		\qty{90}{\kilo\metre\per\hour}
		\item Beschleunigung:
		\qty{9.81}{\metre\per\second\squared}
		\item Kraft: \qty{250}{\newton}
		\item Druck: \qty{1.01325}{\bar}
		\item Energie: \qty{1500}{\joule}
		\item Leistung: \qty{250}{\kilo\watt}
		\item Spannung: \qty{230}{\volt}
		\item Widerstand: \qty{50}{\ohm}
	\end{itemize}
	
	% ==================================================
	\subsection{Einheiten ohne Zahlen mit \texttt{\textbackslash unit}}
	% ==================================================
	
	\begin{itemize}
		\item Länge: 
			\unit{\metre}
		\item Geschwindigkeit:
			\unit{\metre\per\second}
		\item Beschleunigung:
			\unit{\metre\per\second\squared}
		\item Dichte:
			\unit{\kilogram\per\cubic\metre}
		\item Energie:
			\unit{\joule}
		\item Druck:
			\unit{\newton\per\square\metre}
	\end{itemize}
	
	% ==================================================
	\subsection{Einheiten in mathematischen Formeln}
	% ==================================================
	
	\begin{equation}
		F = m \cdot a
	\end{equation}
	
	Beispiel:
	
	\begin{equation}
		F =
		\qty{10}{\kilogram}
		\cdot
		\qty{9.81}{\metre\per\second\squared}
		=
		\qty{98.1}{\newton}
	\end{equation}
	
	Kinetische Energie:
	
	\begin{equation}
		E_{\mathrm{kin}} =
		\frac{1}{2}mv^2
	\end{equation}
	
	Beispiel:
	
	\begin{equation}
		E_{\mathrm{kin}}
		=
		\frac{1}{2}
		\cdot \qty{2}{\kilogram}
		\cdot
		\left(
		\qty{3}{\metre\per\second}
		\right)^2
		=
		\qty{9}{\joule}
	\end{equation}
	
	% ==================================================
	\subsection{Messunsicherheiten}
	% ==================================================
	
	Messwert mit Unsicherheit:
	
	\begin{equation}
		l =
		\qty[
		uncertainty-mode=separate
		]{12.34(5)}{\metre}
	\end{equation}
	
	Temperaturmessung:
	
	\begin{equation}
		T =
		\qty[
		uncertainty-mode=separate
		]{23.5(2)}{\celsius}
	\end{equation}
	
	% ==================================================
	\subsection{Winkel}
	% ==================================================
	
	\begin{itemize}
		\item Winkel: \ang{45}
		\item Dezimalwinkel: \ang{90.5}
		\item Grad, Minuten und Sekunden:
		\ang{30;15;20}
	\end{itemize}
	
	% ==================================================
	\subsection{Bereiche und Listen}
	% ==================================================
	
	\begin{itemize}
		\item Zahlenbereich:
		\numrange{10}{20}
		\item Temperaturbereich:
		\qtyrange{18}{25}{\celsius}
		\item Längenbereich:
		\qtyrange{5}{15}{\metre}
		\item Zahlenliste:
		\numlist{1;2;3;4}
		\item Messwertliste:
		\qtylist{10;20;30}{\kilogram}
	\end{itemize}
	
	% ==================================================
	\subsection{Tabellen mit Dezimalausrichtung}
	% ==================================================
	
	\begin{table}[htbp]
		\centering
		
		\begin{tabular}{
				l
				S[table-format=3.2]
				S[table-format=2.2]
			}
			\toprule
			{Messung} &
			{Masse (\unit{\kilogram})} &
			{Länge (\unit{\metre})} \\
			\midrule
			A & 12.50  & 1.25 \\
			B & 2.35   & 2.50 \\
			C & 125.70 & 12.75 \\
			D & 0.25   & 0.50 \\
			\bottomrule
		\end{tabular}
		
		\caption{Messwerte mit automatisch
			ausgerichteten Dezimalstellen}
		
		\label{tab:siunitx-messwerte}
	\end{table}
	
	% ==================================================
	\subsection{Eigene Einheiten definieren mit \texttt{\textbackslash DeclareSIUnit}}
	% ==================================================
	
	% Beispiel: Energieeinheit
	\DeclareSIUnit{\kWh}{kWh}
	
	Ein Energieverbrauch von
	\qty{125.5}{\kWh}.
	
	% ==================================================
	\subsection{Zusammenfassung}
	% ==================================================
	
	Die wichtigsten Befehle:
	
	\begin{description}
		
		\item[\texttt{\textbackslash num}]
		Zahlen formatieren.
		
		\item[\texttt{\textbackslash unit}]
		Einheiten darstellen.
		
		\item[\texttt{\textbackslash qty}]
		Zahlen mit Einheiten kombinieren.
		
		\item[\texttt{\textbackslash ang}]
		Winkel darstellen.
		
		\item[\texttt{\textbackslash qtyrange}]
		Messwertbereiche darstellen.
		
		\item[\texttt{\textbackslash qtylist}]
		Messwertlisten darstellen.
		
		\item[\texttt{\textbackslash sisetup}]
		Globale Einstellungen definieren.
		
		\item[\texttt{S-Spalte}]
		Dezimalzahlen in Tabellen ausrichten.
		
	\end{description}
	
\end{document}
